\begin{theorem}[Positivity of the rescaling-stable MPDO]%
	\label{thm:mpdo_bnt_label_rescaling_stable_ismpdo}
	\lean{MPOTensor.RescalingStableLengthDependentRFP.mpo_R_entry_formula}
	\lean{MPOTensor.RescalingStableLengthDependentRFP.R_isMPDO}
	\leanok
	\uses{def:mpdo, thm:mpdo_bnt_label_rescaling_stable_length_dependence}
	The operator $\rho^{(N)}(R)$ is positive semidefinite for every
	positive chain length $N$, so $R$ generates MPDOs. Writing each bond matrix
	$R^{pq}$ as a rank-one product scaled by $25/32$, the closed-chain trace
	of a length-$N$ configuration $(\sigma,\tau)$ gives
	\begin{align}
		\rho^{(N)}(R)_{\sigma\tau}
		 & =\left(\frac{25}{32}\right)^N
		\mathbb{1}[\text{cyclic bond match}](\sigma,\tau)
		\prod_n W(\sigma_n,\tau_n),
		\notag
	\end{align}
	where the indicator is one exactly when both $\sigma$ and $\tau$
	satisfy the cyclic bond-matching condition: consecutive letters share
	their bond bit around the ring. In the product, $W$ is evaluated on
	the corresponding bond bits. It is the local factor from
	Theorem~\ref{thm:mpdo_bnt_label_rescaling_stable_length_dependence},
	with eigenvalues $\{1,7/25\}$. Thus $\rho^{(N)}(R)$ is $(25/32)^N$ times
	the Hadamard product of the rank-one bond-matching indicator with the
	pullback of the $N$-fold Kronecker power of $W$.
\end{theorem}
\begin{proof}\leanok
	The entrywise factorization follows by induction on the chain
	length, using the rank-one contraction identity
	\begin{align}
		\tr((u_1v_1^{\mathsf T})\cdots(u_Nv_N^{\mathsf T}))
		 & =(v_N^{\mathsf T}u_1)\prod_{n=1}^{N-1}(v_n^{\mathsf T}u_{n+1}),
		\notag
	\end{align}
	where $u_n,v_n$ are the column vectors in the rank-one factorization
	of the $n$-th bond matrix. Closing the trace produces the cyclic
	bond-matching condition. The bond-matching indicator is positive
	semidefinite because it is an outer product. The pullback of the
	Kronecker power of $W$ is positive semidefinite by congruence.
	The Schur product theorem and the positive scalar $(25/32)^N$
	therefore give positivity.
\end{proof}

\begin{theorem}[Unitality of the transfer map]%
	\label{thm:mpdo_rescaling_transfer_unital}
	\lean{MPOTensor.RescalingStableLengthDependentRFP.transferMap_A_one}
	\leanok
	\uses{thm:mpdo_bnt_label_rescaling_stable_length_dependence}
	For the explicit four-letter tensor $A$, the transfer map
	$\E_A(Y)=\sum_a A^aY(A^a)^\dagger$ satisfies $\E_A(\Id)=\Id$.
\end{theorem}

\begin{proof}\leanok
	Evaluate the four matrix entries of $\E_A(\Id)$ from the
	matrices $A^a$. The diagonal entries are one and the off-diagonal entries
	vanish.
\end{proof}

\begin{theorem}[Traceless diagonal eigenvector of the transfer map]%
	\label{thm:mpdo_rescaling_transfer_diag_alt}
	\lean{MPOTensor.RescalingStableLengthDependentRFP.transferMap_A_diag_alt}
	\leanok
	\uses{thm:mpdo_bnt_label_rescaling_stable_length_dependence}
	For $Z=\diag(1,-1)$, the transfer map satisfies
	$\E_A(Z)=\frac{7}{25}Z$.
\end{theorem}

\begin{proof}\leanok
	Evaluate $\sum_a A^aZ(A^a)^\dagger$ entrywise. Its off-diagonal entries
	vanish, while its diagonal entries are $7/25$ and $-7/25$.
\end{proof}

\begin{theorem}[Annihilation of the first off-diagonal matrix unit]%
	\label{thm:mpdo_rescaling_transfer_single_zero_one}
	\lean{MPOTensor.RescalingStableLengthDependentRFP.transferMap_A_single01}
	\leanok
	\uses{thm:mpdo_bnt_label_rescaling_stable_length_dependence}
	For the matrix unit $E_{01}$, the transfer map satisfies
	$\E_A(E_{01})=0$.
\end{theorem}

\begin{proof}\leanok
	Substitution of the four matrices $A^a$ shows that every entry of
	$\sum_a A^aE_{01}(A^a)^\dagger$ is zero.
\end{proof}

\begin{theorem}[Annihilation of the second off-diagonal matrix unit]%
	\label{thm:mpdo_rescaling_transfer_single_one_zero}
	\lean{MPOTensor.RescalingStableLengthDependentRFP.transferMap_A_single10}
	\leanok
	\uses{thm:mpdo_bnt_label_rescaling_stable_length_dependence}
	For the matrix unit $E_{10}$, the transfer map satisfies
	$\E_A(E_{10})=0$.
\end{theorem}

\begin{proof}\leanok
	Substitution of the four matrices $A^a$ shows that every entry of
	$\sum_a A^aE_{10}(A^a)^\dagger$ is zero.
\end{proof}

\begin{theorem}[Spectra of the local factor and the transfer map]%
	\label{thm:mpdo_bnt_label_rescaling_stable_spectral}
	\lean{MPOTensor.RescalingStableLengthDependentRFP.%
		wMat_eigenvalue_eq_oneLabelChi_entry}
	\lean{MPOTensor.RescalingStableLengthDependentRFP.transferMap_A_eigen}
	\leanok
	\uses{thm:mpdo_bnt_label_rescaling_stable_length_dependence}
	The two diagonal entries of $\chi$ are eigenvalues of the local
	factor $W$: the uniform vector has eigenvalue $1$, and
	the alternating vector has eigenvalue $7/25$. The transfer map
	$\E_A(Y)=\sum_a A^aY(A^a)^\dagger$ has eigenvalues $1,7/25,0,0$,
	counted with multiplicity.
	The identity is fixed, the traceless diagonal matrix
	$\diag(1,-1)$ has eigenvalue $7/25$, and the
	off-diagonal matrix units are annihilated. Thus the entries of $\chi$
	are also the nonzero eigenvalues of $\E_A$.
	% The explicit one-label tensor-attached clause below uses these two
	% eigenvalues as the diagonal entries of its chi matrix.
\end{theorem}
\begin{proof}\leanok
	\uses{thm:mpdo_rescaling_transfer_unital, thm:mpdo_rescaling_transfer_diag_alt, thm:mpdo_rescaling_transfer_single_zero_one, thm:mpdo_rescaling_transfer_single_one_zero}
	The uniform and alternating vectors are eigenvectors of the explicit
	$2\times2$ matrix $W$:
	\begin{align}
		W\begin{pmatrix}1\\1\end{pmatrix}
		 & =\begin{pmatrix}1\\1\end{pmatrix},
		\notag\\
		W\begin{pmatrix}1\\-1\end{pmatrix}
		 & =\frac{7}{25}\begin{pmatrix}1\\-1\end{pmatrix}.
		\notag
	\end{align}
	Direct expansion of the letter matrices gives the action on a basis
	of $\MN{2}$:
	\begin{align}
		\E_A(\Id) & =\Id,
		\notag\\
		\E_A(\diag(1,-1)) & =\frac{7}{25}\diag(1,-1),
		\notag\\
		\E_A(E_{01}) & =0,
		\notag\\
		\E_A(E_{10}) & =0.
		\notag
	\end{align}
\end{proof}

\begin{lemma}[Injectivity of \texorpdfstring{$R^{\bullet\bullet}$}{the doubled-index tensor}]%
	\label{lem:mpdo_bnt_label_rescaling_stable_injective}
	\lean{MPOTensor.RescalingStableLengthDependentRFP.R_toMPSTensor_isInjective}
	\leanok
	\uses{def:injective, thm:mpdo_bnt_label_rescaling_stable_length_dependence}
	The $16$ letters $R^{pq}$ of the doubled-index tensor
	$R^{\bullet\bullet}$ span the full bond matrix algebra $\MN{4}$.
\end{lemma}

\begin{proof}\leanok
	Each letter $R^{pq}$ is, up to a nonzero scalar, one of the $16$
	matrix units of $\MN{4}$. Indeed, the outer-product
	structure of $R^{pq}$ gives a single matrix unit because
	each factor $A^a$, $A^b$ of the vertical letters
	$B^{ab}=A^a\otimes\overline{A^b}$ is a scaled matrix unit of $\MN{2}$.
	Hence the standard matrix-unit basis lies in the span of the letters.
\end{proof}

\begin{theorem}[Canonical form of \texorpdfstring{$R$}{R}]%
	\label{thm:mpdo_bnt_label_rescaling_stable_cf}
	\lean{MPOTensor.RescalingStableLengthDependentRFP.R_toMPSTensor_isCPSVCanonicalForm}
	\leanok
	\uses{def:cpsv_canonical_form, def:nt_cpsv, thm:mpdo_bnt_label_rescaling_stable_length_dependence}
	The doubled-index tensor $R^{\bullet\bullet}$, with bond dimension $4$
	and $16$ letters $R^{pq}$ for $p,q\in\{0,\ldots,3\}$, is in CF as in
	Definition~\ref{def:cpsv_canonical_form}
	(\ref{eq:cpsv_canonical_form}).
	It consists of a single normal block occupying the full bond
	space ($r=1$, $D_1=4$, $U=\Id$), with weight
	$\mu=\sqrt{337/512}$. This is the unique positive weight for which
	the normalized block $\mu^{-1}R^{\bullet\bullet}$ has an idempotent
	transfer map, and hence spectral radius one.
\end{theorem}
\begin{proof}\leanok
	\uses{lem:mpdo_bnt_label_rescaling_stable_injective}
	By Lemma~\ref{lem:mpdo_bnt_label_rescaling_stable_injective}, the $16$
	letters span the full bond matrix algebra. Thus
	$R^{\bullet\bullet}$ is injective, and no nontrivial invariant orthogonal
	projection exists.

	The transfer map
	$\E_{R^{\bullet\bullet}}(X)=\sum_{p,q}R^{pq}X(R^{pq})^\dagger$
	has rank one: every image is a scalar multiple of a fixed diagonal
	matrix, and the scalar evaluated at that matrix is $337/512$.
	Rescaling the tensor by $\mu^{-1}=\sqrt{512/337}$ makes its transfer map
	idempotent. Together with injectivity, this gives a primitive transfer
	map with spectral radius one; all its other eigenvalues are zero.
	Thus $\mu^{-1}R^{\bullet\bullet}$ is an NT with the required
	normalization. The equality
	$R^{\bullet\bullet}=\mu(\mu^{-1}R^{\bullet\bullet})$ gives exact
	reconstruction with the identity coisometry.
\end{proof}

\begin{theorem}[Non-nilpotency of the retained horizontal block]%
	\label{thm:mpdo_bnt_label_rescaling_stable_non_nilpotent}
	\lean{MPOTensor.RescalingStableLengthDependentRFP.%
		physTraceTransfer_R_not_isNilpotent}
	\lean{MPOTensor.RescalingStableLengthDependentRFP.%
		doubledPhysTraceTransfer_retainedBlock}
	\lean{MPOTensor.RescalingStableLengthDependentRFP.%
		doubledPhysTraceTransfer_retainedBlock_not_isNilpotent}
	\leanok
	\uses{def:mpdo_doubled_phys_trace_transfer, def:mpdo_simple_tensor, thm:mpdo_bnt_label_rescaling_stable_length_dependence, thm:mpdo_bnt_label_rescaling_stable_cf}
	Let $\widehat R=\mu^{-1}R^{\bullet\bullet}$ be the single normal
	block in this CF, where $\mu=\sqrt{337/512}$. Then
	\begin{align}
		\mathcal T_{\widehat R} & =\mu^{-1}\mathcal T_R,
		\notag\\
		\mathcal T_R^2 & =\mathcal T_R,
		\notag\\
		(\mathcal T_R)_{00} & =\frac12.
		\notag
	\end{align}
	Consequently neither $\mathcal T_R$ nor
	$\mathcal T_{\widehat R}$ is nilpotent. Thus the retained BNT element
	satisfies the non-nilpotency condition in the definition of simplicity.

	This conclusion concerns that condition for this one-block
	presentation only, not other presentations or blockings.
\end{theorem}
\begin{proof}\leanok
	\uses{def:mpdo_doubled_phys_trace_transfer, thm:mpdo_bnt_label_rescaling_stable_length_dependence, thm:mpdo_bnt_label_rescaling_stable_cf}
	A nilpotent idempotent is zero, but the $(0,0)$ entry of
	$\mathcal T_R$ is
	\begin{align}
		\frac{25}{32}\tr(A^0)^2
		& =\frac{25}{32}\left(\frac45\right)^2
		=\frac12,
		\notag
	\end{align}
	so $\mathcal T_R$ is not nilpotent. The diagonal physical contraction is
	linear in the tensor, which gives
	$\mathcal T_{\widehat R}=\mu^{-1}\mathcal T_R$. Since $\mu>0$, a
	nilpotent $\mathcal T_{\widehat R}$ would remain nilpotent after
	multiplication by $\mu$, forcing $\mathcal T_R$ to be nilpotent.
\end{proof}

\begin{theorem}[Simplicity of the rescaling-stable dimer]
	\label{thm:mpdo_bnt_label_rescaling_stable_simple}
	\lean{MPOTensor.RescalingStableLengthDependentRFP.oneSiteDoubledEquiv}
	\lean{MPOTensor.RescalingStableLengthDependentRFP.%
		oneSiteDoubledEquiv_diagonal}
	\lean{MPOTensor.RescalingStableLengthDependentRFP.%
		toMPSTensor_blockTensor_R_one}
	\lean{MPOTensor.RescalingStableLengthDependentRFP.R_isSimple}
	\leanok
	\uses{def:mpdo_simple_tensor, thm:mpdo_bnt_label_rescaling_stable_length_dependence}
	The dimer tensor $R$ is simple, with blocking length $L=1$.
	After the canonical relabeling of the doubled physical alphabet,
	the BNT has one member, $\widehat R$, and its physical-trace
	transfer is non-nilpotent.
\end{theorem}

\begin{proof}\leanok
	\uses{thm:cpsv_representative_sector_decomposition, thm:mpdo_bnt_label_rescaling_stable_ismpdo, thm:mpdo_bnt_label_rescaling_stable_cf, thm:mpdo_bnt_label_rescaling_stable_non_nilpotent}
	Theorem~\ref{thm:mpdo_bnt_label_rescaling_stable_ismpdo} gives the MPDO
	condition. Choose the one-site blocking. Its doubled-index tensor is the
	canonical physical relabeling of $R^{\bullet\bullet}$. Refine the CF
	from Theorem~\ref{thm:mpdo_bnt_label_rescaling_stable_cf} to its
	BNT sector presentation. The only representative is $\widehat R$, whose
	physical-trace transfer is non-nilpotent by
	Theorem~\ref{thm:mpdo_bnt_label_rescaling_stable_non_nilpotent}.
\end{proof}

\begin{theorem}[The rescaling-stable tensor is a renormalization fixed
		point]%
	\label{thm:mpdo_bnt_label_rescaling_stable_is_rfp_via_ts}
	\lean{MPOTensor.RescalingStableLengthDependentRFP.isRFPViaTS_R}
	\leanok
	\uses{def:is_rfp_via_ts, thm:mpdo_bnt_label_rescaling_stable_length_dependence}
	The tensor $R$ of
	Theorem~\ref{thm:mpdo_bnt_label_rescaling_stable_length_dependence} is an
	RFP in the sense of Definition~\ref{def:is_rfp_via_ts}.
	The tpCPM $T$ from one site to two sites has
	two Kraus operators built from the eigenvectors of $W$, with amplitude
	$5/8$ on the uniform eigenvector and $\sqrt7/8$ on the alternating
	eigenvector. The tpCPM $S$ from two sites to one site has four Kraus
	operators. Two use the amplitude $1/\sqrt2$ and the eigenvectors of
	$W$ on physical-index pairs satisfying the bond-matching condition.
	The other two assign the pairs that fail this condition deterministically,
	so that the four Kraus operators resolve the identity.
	These maps satisfy $T[M_1(X)]=M_2(X)$ and $S[M_2(X)]=M_1(X)$
	for every boundary matrix $X$, with $M=R$.
\end{theorem}
\begin{proof}\leanok
	Identify each physical index $p\in\{0,1,2,3\}$ with a pair of bits
	$(p',p'')\in\{0,1\}^2$, and write $i_1\ast i_2:=(i_1',i_2'')$
	for the physical index obtained from the first bit of $i_1$ and the second
	bit of $i_2$. For $M=R$, the two-site physical operator satisfies
	\begin{align}
		M_2(X)_{(i_1,i_2)(j_1,j_2)}
		 & =\frac{25}{32}
		\mathbb{1}[i_1''=i_2']\,\mathbb{1}[j_1''=j_2']
		W(i_1'',j_1'')\,M_1(X)_{i_1\ast i_2,\,j_1\ast j_2},
		\notag
	\end{align}
	and therefore vanishes whenever either bond match fails.
	Write $e_0=(1,1)$ and $e_1=(1,-1)$ for the eigenvectors of $W$,
	with eigenvalues $\lambda_0=1$ and $\lambda_1=7/25$
	(Theorem~\ref{thm:mpdo_bnt_label_rescaling_stable_spectral}).
	The two Kraus operators of $T$ have amplitudes
	$\sqrt{25\lambda_s/64}$ on $e_s$. For $b,b'\in\{0,1\}$,
	the eigendecomposition of $W$ gives
	\begin{align}
		\sum_{s\in\{0,1\}}\frac{25}{64}\lambda_s e_s(b)e_s(b')
		 & =\frac{25}{32}W(b,b').
		\notag
	\end{align}
	This reproduces the coefficient in $M_2(X)$, so $T[M_1(X)]=M_2(X)$.

	The two Kraus operators of $S$ supported on bond-matching pairs have
	amplitude $1/\sqrt2$ on $e_0,e_1$. For $b,b'\in\{0,1\}$,
	\begin{align}
		\frac12\sum_{s\in\{0,1\}}e_s(b)e_s(b')
		 & =\begin{cases}1,&b=b',\\0,&b\ne b'.\end{cases}
		\notag
	\end{align}
	The diagonal case gives the resolution of the identity on the
	bond-matching subspace, and the off-diagonal case gives cancellation
	between these two Kraus operators. The other two Kraus operators
	assign the bond-mismatched pairs deterministically.
	They do not contribute to $S[M_2(X)]$, because $M_2(X)$ vanishes
	on these pairs. Together the four Kraus operators resolve the identity
	on the full two-site domain and give $S[M_2(X)]=M_1(X)$.
\end{proof}

% Source: arXiv:1606.00608, Theorem 4.14(ii), lines 972--985,
% and Appendix C.4, lines 1929--1947 and 2011--2029.
\begin{theorem}[Vertical BNT and algebra structure for \texorpdfstring{$R$}{R}]%
	\label{thm:mpdo_bnt_label_rescaling_stable_bnt_algebra_clause}
	\lean{MPOTensor.RescalingStableLengthDependentRFP.%
		R_oneLabelBNTAlgebraTensorClause}
	\leanok
	\uses{def:mpdo_bnt_algebra_tensor_clause, thm:mpdo_bnt_label_rescaling_stable_length_dependence}
	The tensor $R$ has a vertical BNT with one element $M_0$.
	This component has $16$ physical letters, bond dimension $4$, and
	\begin{align}
		M_0^{(a,b)} & =A^a\otimes\overline{A^b}.
		\notag
	\end{align}
	The positive diagonal multiplicity and coefficient matrices are
	\begin{align}
		\mu_0 & =\begin{pmatrix}\frac{25}{32}\end{pmatrix},
		\notag\\
		\chi_{0,0,0} & =\diag\!\left(1,\frac{7}{25}\right).
		\notag
	\end{align}
	The vertical tensor is reconstructed exactly as
	\begin{align}
		\operatorname{vertical}(R) & =\frac{25}{32}M_0.
		\notag
	\end{align}
	For every $L>0$, the vertical operator satisfies
	\begin{align}
		O_L(M_0)^2
		 & =\left(1+\left(\frac{7}{25}\right)^L\right)O_L(M_0)
		=\tr\!\left(\chi_{0,0,0}^L\right)O_L(M_0).
		\notag
	\end{align}
	At length one, $m_0=\tr(\mu_0)=25/32$ and
	$c^{(1)}_{0,0,0}=32/25$. Thus the vector $m=(m_0)$ is an idempotent
	for the multiplication induced by $c^{(1)}$, since
	$m_0=c^{(1)}_{0,0,0}m_0^2$:
	\begin{align}
		\frac{25}{32}
		 & =\frac{32}{25}\left(\frac{25}{32}\right)^2.
		\notag
	\end{align}
\end{theorem}

\begin{proof}\leanok
	\uses{thm:is_normal_tensor_of_is_normal_left_canonical, lem:mpo_doubled, thm:mpdo_bnt_label_rescaling_stable_length_dependence, thm:mpdo_bnt_label_rescaling_stable_chi_uniform}
	Each letter of $M_0$ is a nonzero multiple of a matrix unit, so the letters
	span $\MN{4}$. The tensor $A$ is left-canonical; hence $M_0$ is
	left-canonical and algebraically normal, and therefore is an NT.
	The identity coisometry gives both the vertical decomposition and
	its exact reconstruction.

	The vertical operator of $M_0=A\otimes\overline A$ is
	$\ketbra{V^{(L)}(A)}{V^{(L)}(A)}$. Its square is its self-overlap times
	itself. A direct transfer-map computation gives
	$\braket{V^{(L)}(A)}{V^{(L)}(A)}=\tr(W^L)$, while
	Walsh--Hadamard diagonalization gives
	$\tr(W^L)=\tr(\chi_{0,0,0}^L)=1+(7/25)^L$.
	Finally, $c^{(1)}_{0,0,0}m_0^2=(32/25)(25/32)^2=m_0$.
\end{proof}

\begin{remark}[Vertical and horizontal contractions]
	Let $\operatorname{rot}(\operatorname{vertical}(R))$ denote the MPO obtained
	by separating the $16$-valued physical index of the vertical tensor into
	two four-valued indices. Then
	\begin{align}
		\rho^{(L)} \left(\operatorname{rot}(\operatorname{vertical}(R))\right)
		 & =\left(\frac{25}{32}\right)^L O_L(M_0).
		\notag
	\end{align}
	This is a rotated vertical contraction, not the original horizontal tensor
	$R$, even though both physical and bond dimensions equal $4$. At length one,
	$\rho^{(1)}(R)$ has rank two, whereas $O_1(M_0)$ has rank one, so these
	operators are not proportional. In particular, no identity
	$\rho^{(L)}(R)=(25/32)^L O_L(M_0)$ is asserted.
\end{remark}

\begin{theorem}[Closed-operator factorization and local diagonalization for \texorpdfstring{$R$}{R}]%
	\label{thm:mpdo_bnt_label_rescaling_stable_chi_attachment}
	\lean{MPOTensor.RescalingStableLengthDependentRFP.%
		mpo_R_eq_B_mul_wN_mul_transpose}
	\lean{MPOTensor.RescalingStableLengthDependentRFP.B_transpose_mul_B}
	\lean{MPOTensor.RescalingStableLengthDependentRFP.%
		wMat_eq_conj_diagonal_oneLabelChi}
	\lean{MPOTensor.RescalingStableLengthDependentRFP.%
		hadamard2_mul_diagonal_oneLabelChi_mul_hadamard2}
	\leanok
	\uses{thm:mpdo_bnt_label_rescaling_stable_length_dependence, thm:mpdo_bnt_label_rescaling_stable_ismpdo}
	For every positive chain length $N$, the horizontal operator factors as
	\begin{align}
		\rho^{(N)}(R) & =\left(\frac{25}{32}\right)^N B\,W_N\,B^{\mathsf T},
		\notag
	\end{align}
	where $B$ sends each bit-string basis vector to the basis vector of the
	unique cyclic-bond-matching physical-index string with those first bits,
	and $B^{\mathsf T}B=\Id$. Here
	$W_N(a,b)=\prod_n W(a_n,b_n)$ is the $N$-fold Kronecker power of $W$.
	For the Walsh--Hadamard matrix
	$H=\begin{pmatrix}1&1\\1&-1\end{pmatrix}$ and
	$\chi=\diag(1,\tfrac{7}{25})$, one has $H\chi H=2W$.
	Thus the entries of $\chi$ are the eigenvalues of the local factor $W$.
	Theorem~\ref{thm:mpdo_bnt_label_rescaling_stable_bnt_algebra_clause}
	identifies this same matrix as the coefficient matrix
	$\chi_{0,0,0}$ of a vertical BNT with one element.
	The factorization here concerns the original horizontal operator
	$\rho^{(N)}(R)$, not the rotated vertical contraction.
\end{theorem}
\begin{proof}\leanok
	Let $\phi$ read the first bits of a cyclic-bond-matching string.
	Given a bit string $a$, each adjacent pair of entries of $a$ determines
	one physical index around the ring. This gives the unique matching
	string whose image under $\phi$ is $a$, and defines $B$.
	Injectivity gives $B^{\mathsf T}B=\Id$. Rewriting the entrywise formula
	of Theorem~\ref{thm:mpdo_bnt_label_rescaling_stable_ismpdo} through this
	embedding gives
	$\rho^{(N)}(R)=(25/32)^N B W_N B^{\mathsf T}$.
	The identity $H\chi H=2W$ follows by direct $2\times2$ computation.
\end{proof}

\begin{theorem}[Trace-power formula for the structure coefficients]%
	\label{thm:mpdo_bnt_label_rescaling_stable_chi_uniform}
	\lean{MPOTensor.RescalingStableLengthDependentRFP.oneLabelChiTracePowerForm}
	\lean{MPOTensor.RescalingStableLengthDependentRFP.%
		wMat_pow_trace_eq_oneLabelChi_matrix_pow_trace}
	\lean{MPOTensor.RescalingStableLengthDependentRFP.%
		oneLabelCoeffs_coeff_eq_wMat_pow_trace}
	\leanok
	\uses{def:mpdo_positive_bnt_label_chi_trace_power_form, lem:mpdo_positive_bnt_label_chi_eq_trace_matrix_pow, thm:mpdo_bnt_label_rescaling_stable_length_dependence, thm:mpdo_bnt_label_rescaling_stable_chi_attachment}
	The coefficients $c^{(L)}_{0,0,0}$ have a trace-power representation
	with the fixed positive diagonal matrix
	$\chi=\diag(1,\tfrac{7}{25})$, as in
	Definition~\ref{def:mpdo_positive_bnt_label_chi_trace_power_form}.
	For every $L>0$,
	\begin{align}
		c^{(L)}_{0,0,0} & =\tr(\chi^L)=\tr(W^L)=1+(7/25)^L.
		\notag
	\end{align}
	The equality of traces follows from the diagonalization in
	Theorem~\ref{thm:mpdo_bnt_label_rescaling_stable_chi_attachment}.
	Here $W^L$ is the ordinary matrix power of the fixed $2\times2$
	matrix $W$, not the Kronecker power $W_L$ in $\rho^{(L)}(R)$.
	Theorem~\ref{thm:mpdo_bnt_label_rescaling_stable_bnt_algebra_clause}
	identifies these coefficients with the structure coefficients of the
	vertical BNT $\{M_0\}$. It does not identify $O_L(M_0)$ with the
	original horizontal operator $\rho^{(L)}(R)$.
\end{theorem}
\begin{proof}\leanok
	The coefficients are defined by the diagonal entries of $\chi$.
	Since these entries are positive,
	Lemma~\ref{lem:mpdo_positive_bnt_label_chi_eq_trace_matrix_pow}
	gives $c^{(L)}_{0,0,0}=\tr(\chi^L)$ for every $L>0$.
	The Walsh--Hadamard identity $H\chi H=2W$ and $HH=2\Id$ give
	$H\chi=WH$. Induction then gives $H\chi^L=W^L H$.
	Conjugating by $H$, whose inverse is $H/2$, and using cyclicity
	of the trace yields $\tr(W^L)=\tr(\chi^L)$ for every $L$.
\end{proof}
